\documentclass[11pt]{article}
\usepackage[a4paper,margin=2.35cm]{geometry}
\usepackage{fontspec}
\usepackage{amsmath,amssymb}
\setmainfont[ItalicFont=Noto Serif Italic]{Malgun Gothic}
\setsansfont{Malgun Gothic}
\setlength{\parindent}{1.2em}
\setlength{\parskip}{0.45em}
\setlength{\emergencystretch}{3em}
\providecommand{\srcspaced}[1]{\textbf{#1}}
\begin{document}

더 나아가 다음이 성립한다. \srcspaced{유리수체에 관하여 차수 $2^n$인 순환체 $\Omega_n$ 가운데 $(-1)$을 두 제곱의 합으로 나타낼 수 있는 것들은 모두 사원수체의 극소 분해체이다.} 실제로 $(-1)$을 제곱들의 합으로 나타낼 수 있으므로 분해체는 실수체일 수 없다. 한편 $\Omega_n$의 차수 $2^{n-1}$인 부분체---따라서 $\Omega_n$의 모든 진부분체---는 반드시 실수체이다. 이 부분체는 $\Omega_n$과 모든 실대수적 수의 체의 교집합이기 때문이다. 실제로 $\Omega_n$은 이 교집합에 대한 차수 2의 갈루아체가 되므로, 이 교집합은 반드시 차수 $2^{n-1}$인 그 부분체와 일치한다.

\end{document}
